\subsection{Systemprüfung und Installation}
\label{subsec:systempruefung-installation}

Der nur lesende \texttt{doctor} prüft profilabhängig Laufzeiten, Konfiguration,
Pfade, Abhängigkeiten und Ports. Deaktivierte Schichten sind kein Fehler.
Fehlendes Docker erzeugt allgemein nur eine Warnung; Container-, Datenbank- und
Tauri-Voraussetzungen besitzen strengere eigene Diagnosewege
\parencite{kleiveist2026tooling,kleiveist2026configuration}.

\texttt{install} richtet nur aktive Komponenten ein: Frontend-Abhängigkeiten
werden bevorzugt aus dem Lockfile installiert, Backend-Abhängigkeiten in einer
virtuellen Umgebung. Die dedizierte \path{tools/.venv} wird profilunabhängig
für Repository-Werkzeuge einschließlich Ruff eingerichtet; die
\path{backend/.venv} bleibt davon getrennt. Datenbankpakete und Chromium kommen
nur bei Bedarf hinzu; für Python besteht neben \texttt{uv} ein
pip/venv-Fallback. Eine lokale \texttt{.env} bleibt unverändert. Native
Betriebssystem- und Rust-Voraussetzungen behandelt \texttt{tauri install}
separat.
Konkret verwendet das Frontend bei vorhandenem Lockfile \texttt{npm ci}. Ein
aktives Backend erhält eine eigene Umgebung und nur die profilabhängigen
Anforderungssätze. Chromium wird ausschließlich bei konfiguriertem Playwright
installiert; deaktivierte Schichten werden übersprungen.

Am historischen Beobachtungsstand \texttt{461fc751} bestand für generierte
Backendprofile eine für Governance relevante Abweichung der Installation:
\texttt{install} konnte Ruff über \path{backend/.venv} bereitstellen und auf
\path{tools/.venv} verzichten, während der Quality-Adapter Ruff dort oder im
globalen Suchpfad erwartete. In einem frisch erzeugten
\texttt{web-cloud}-Profil endete die Installation erfolgreich, der
anschließende Quality-Befehl meldete Ruff jedoch als nicht verfügbar. Der
Core-CI-Job vermied diese Konstellation durch
\texttt{install --skip-backend}; die generierten Backend-Profiljobs taten dies
nicht. Diese commitgebundene Beobachtung erklärt die damaligen Profilfehler
und bleibt Teil der historischen Evidenz
\parencite{kleiveist2026quality-implementation,kleiveist2026quality-ci}.

Im v1.0.0-Kandidatenstand ist die Abweichung behoben: \texttt{install} stellt
\path{tools/.venv} unabhängig vom Profil bereit, und der Quality-Adapter nutzt
diese dedizierte Tooling-Runtime. Der lokale Pfad \texttt{install}
$\rightarrow$ \texttt{quality} besitzt damit hinsichtlich der Ruff-Auflösung
einen einheitlichen Vertrag. Diese spätere Korrektur schreibt die
commitgebundenen CI-Ergebnisse der Baseline nicht um.

Die README fordert Git, Python ab 3.11, Node.js ab 24 und für Desktopprofile
Rust Stable. Der allgemeine Doctor prüft Git und diese Versionsbereiche nicht;
der Tauri-Doctor akzeptiert jede aktive Rust-Toolchain. CI richtet die
festgelegten Python- und Node-Hauptversionen dagegen explizit ein. Damit weicht
die lokale Durchsetzung von der Dokumentation ab
\parencite{kleiveist2026governance-snapshot,kleiveist2026ci}.
Diese Grenze betrifft die Prüfung der Voraussetzung, nicht deren dokumentierte
Gültigkeit.
\beginnerinput{workspace/statement/500_tooling_workflow/530}
